% =====================================================================
%  PRIMA DI INIZIARE: la dispensa e l'Allenatore (capitolo non numerato,
%  così i moduli restano Capitolo 1..14 e gli esercizi N.k non cambiano)
% =====================================================================
\chapter*{Prima di iniziare: dispensa e Allenatore}
\addcontentsline{toc}{chapter}{Prima di iniziare: dispensa e Allenatore}
\markboth{Prima di iniziare}{Prima di iniziare}

Questo materiale è un \textbf{kit} pensato per preparare l'esame di Programmazione 2
dall'inizio alla fine. È fatto di due parti che lavorano insieme:
\begin{itemize}
  \item la \textbf{dispensa} che stai leggendo, con il volume separato delle
        \textbf{Soluzioni}: è dove si \emph{capisce} (teoria, disegni della memoria,
        codice commentato);
  \item l'\textbf{Allenatore}, un'app per Windows con la versione per telefono (Android e
        iPhone): è dove si \emph{allena} quello che si è capito. Legge i quiz, le
        domande d'orale e gli esercizi \textbf{direttamente da questa dispensa}, compila il
        tuo codice con \cc|g++| come in laboratorio, tiene traccia dei progressi e ti
        accompagna con un piccolo compagno, \textbf{Bit}, che cresce con te.
\end{itemize}
Leggere senza allenarsi dà l'illusione di sapere; allenarsi senza aver capito è
frustrante. Questo capitolo spiega come installare tutto e come usare le due parti
insieme, giorno per giorno.

% ---------------------------------------------------------------------
\section*{Cosa c'è nel kit}
\addcontentsline{toc}{section}{Cosa c'è nel kit}
% ---------------------------------------------------------------------

\begin{center}
\begin{tabular}{p{4.6cm}p{9.6cm}}
\toprule
\textbf{File o cartella} & \textbf{A cosa serve} \\
\midrule
\texttt{dispensa-light.pdf} & la dispensa su sfondo bianco, da stampare \\
\texttt{dispensa-dark.pdf} & la stessa su sfondo scuro, da leggere a schermo \\
\texttt{soluzioni-*.pdf} & le soluzioni (light e dark) dei 10 esercizi di ogni modulo e della simulazione \\
\texttt{Dispensa/codice/} & tutti i programmi della dispensa, divisi per modulo (\texttt{m01}\ldots\texttt{m14}), con input e output reali \\
\texttt{Allenatore Prog2.exe} & l'app per Windows (si scarica dalla pagina del progetto) \\
versione per telefono & un indirizzo web che si apre dal telefono e si installa come un'app \\
\bottomrule
\end{tabular}
\end{center}

% ---------------------------------------------------------------------
\section*{Installare l'Allenatore su Windows}
\addcontentsline{toc}{section}{Installare l'Allenatore su Windows}
% ---------------------------------------------------------------------

\begin{passopasso}[: installazione]
\begin{enumerate}
  \item Scarica \texttt{Allenatore Prog2.exe} dalle \emph{Releases} del repository
        \url{https://github.com/KurtSixxxx/MyUniversity/releases} e mettilo \textbf{accanto alla cartella} \texttt{Dispensa}:
        \begin{center}\ttfamily
        \begin{tabular}{l}
        Prog2/ \\
        \quad Allenatore Prog2.exe \\
        \quad Dispensa/ \quad (con modulo01.tex \ldots\ e codice/)
        \end{tabular}
        \end{center}
        Così l'app legge la dispensa \q{dal vivo}. Se la cartella non c'è, usa una copia
        dei contenuti inclusa nell'eseguibile: funziona lo stesso.
  \item Fai doppio clic. La prima volta Windows può mostrare \q{Windows ha protetto il
        PC} (SmartScreen): è normale per i programmi non firmati. Scegli \textbf{Ulteriori
        informazioni} $\to$ \textbf{Esegui comunque}.
  \item Al primo avvio l'app chiede nome, data dell'appello, docente, colore e il nome
        del tuo compagno. Tutto si cambia quando vuoi in \textbf{Impostazioni}
        (\texttt{Ctrl+,}).
  \item Installa il compilatore (passo successivo): senza \cc|g++| quiz e orale
        funzionano, ma gli esercizi non si possono eseguire.
\end{enumerate}
\end{passopasso}

\subsection*{Il compilatore g++ (MSYS2)}

All'esame si compila con \cc|g++|. Su Windows il modo più semplice per averlo è MSYS2.
\begin{passopasso}[: installare g++]
\begin{enumerate}
  \item Scarica l'installatore da \texttt{msys2.org} e installalo nella cartella proposta
        (\cc|C:\msys64|).
  \item Apri dal menu Start il terminale \textbf{MSYS2 UCRT64} e scrivi:
\begin{lstlisting}[language=bash]
pacman -S mingw-w64-ucrt-x86_64-gcc
\end{lstlisting}
        Conferma con Invio quando chiede.
  \item Aggiungi \cc|C:\msys64\ucrt64\bin| alla variabile di sistema \texttt{Path}:
        cerca \q{variabili d'ambiente} nel menu Start $\to$ \emph{Variabili d'ambiente}
        $\to$ \texttt{Path} $\to$ \emph{Modifica} $\to$ \emph{Nuovo}.
  \item Apri un nuovo Prompt dei comandi e verifica con \cc|g++ --version|. Poi
        riavvia l'Allenatore: in Impostazioni $\to$ Sistema vedrai \q{g++ trovato}.
\end{enumerate}
\end{passopasso}

\begin{consiglio}[: VS Code, facoltativo]
L'editor dell'app basta per tutti gli esercizi. Se preferisci VS Code (all'esame potresti
trovarlo), in Impostazioni $\to$ Esercizi in VS Code l'app crea una cartella con un file
\texttt{.cpp} per ogni esercizio, con la consegna già scritta in commento.
\texttt{Ctrl+Shift+B} compila ed esegue; poi nell'app premi \textbf{Importa da VS Code}
per confrontare l'output con la soluzione.
\end{consiglio}

% ---------------------------------------------------------------------
\section*{Le cinque schermate}
\addcontentsline{toc}{section}{Le cinque schermate}
% ---------------------------------------------------------------------

\begin{description}
  \item[Oggi (\texttt{Ctrl+1}).] Il conto alla rovescia verso l'appello, il
        \textbf{giorno del piano di studio} (Appendice \q{Verso l'esame}) con i pulsanti
        \emph{Quiz di oggi} ed \emph{Esercizi di oggi}, le missioni del giorno, i tuoi
        progressi per modulo e i \textbf{consigli}: l'app guarda dove sbagli e ti propone
        cosa rifare (le risposte sbagliate, il modulo più debole, il prossimo
        laboratorio).
  \item[Quiz (\texttt{Ctrl+2}).] Le domande a crocette di tutti i moduli, prese dalla
        sezione 5 di ogni capitolo, con lo stesso commento della soluzione. Quattro
        modalità: \emph{Allenamento} (correzione subito; le domande sbagliate o mai viste
        escono più spesso), \emph{Da correggere}, \emph{Mai viste} ed \emph{Esame}
        (30 domande in 60 minuti, correzione solo alla consegna, soglia 18). Si risponde
        anche da tastiera: \texttt{A}--\texttt{D}, Invio per andare avanti.
  \item[Orale (\texttt{Ctrl+3}).] Le domande della sezione 7 come \textbf{flashcard}.
        Rispondi \textbf{ad alta voce} (c'è un cronometro da 90 secondi), poi guarda la
        risposta modello e valutati: \emph{Non la sapevo}, \emph{In parte}, \emph{La
        sapevo}. Ogni domanda sale di livello (da 0 a 4) quando la sai: quelle a livello
        basso tornano più spesso. È la ripetizione distanziata dell'appendice sul metodo.
  \item[Codice (\texttt{Ctrl+4}).] Per ogni modulo: il \emph{pattern} da riscrivere a
        memoria, il \emph{laboratorio a step} e i 10 esercizi, più la simulazione
        d'esame. Scrivi nell'editor, premi \texttt{Ctrl+Invio}: l'app compila con
        \cc|g++ -Wall -Wextra| come in laboratorio, segna in rosso le righe con errori e
        mostra l'output. \emph{Confronta} lo mette a fianco dell'output della soluzione;
        \emph{Soluzione} la mostra (dopo una conferma: prima provaci davvero).
  \item[Bit (\texttt{Ctrl+5}).] Il tuo compagno di studio, il livello, le missioni, i
        traguardi e da dove vengono i tuoi punti (sezione successiva).
\end{description}

% ---------------------------------------------------------------------
\section*{Come usare dispensa e Allenatore insieme}
\addcontentsline{toc}{section}{Come usare dispensa e Allenatore insieme}
% ---------------------------------------------------------------------

Ogni modulo della dispensa ha le stesse otto sezioni, e ognuna ha il suo posto nell'app.
La regola è sempre la stessa: \textbf{prima capisci sulla dispensa, poi ti alleni
nell'app, poi torni sulla dispensa dove l'app ti dice che sbagli}.

\begin{center}\small
\begin{tabular}{p{4.3cm}p{4.2cm}p{5.6cm}}
\toprule
\textbf{Sezione del modulo} & \textbf{Dove la alleni} & \textbf{Come} \\
\midrule
1. Teoria \& analogia & dispensa, con carta e penna & leggi, poi chiudi e riscrivi l'idea a parole tue \\
2. Disegni della memoria & dispensa & ridisegnali a mano, senza guardare \\
3. Pattern di codice & Codice $\to$ Pattern & riscrivilo a memoria finché compila e l'output coincide \\
4. Trappole d'esame & dispensa & rileggile prima dei quiz del modulo \\
5. Quiz a crocette & Quiz $\to$ Allenamento & solo il modulo; rifai le sbagliate il giorno dopo \\
6. Laboratorio a step & Codice $\to$ Laboratorio & uno step alla volta, compilando sempre \\
7. Domande dell'orale & Orale & ad alta voce, con il cronometro \\
8. Esercizi al PC & Codice $\to$ Esercizi & in ordine di difficoltà; Confronta, poi Soluzione \\
\bottomrule
\end{tabular}
\end{center}

\subsection*{Una giornata tipo}

\begin{passopasso}[: dispensa + app in una giornata]
\begin{enumerate}
  \item Apri l'app su \textbf{Oggi}: ti dice l'argomento del piano (se hai impostato la
        data dell'appello) e le tre missioni del giorno.
  \item \textbf{Teoria} sulla dispensa (sezioni 1--4 del modulo del giorno).
  \item \textbf{Quiz di oggi} dall'app: 10--20 domande del modulo. Per ogni risposta
        sbagliata rileggi il commento e, se serve, il paragrafo della dispensa.
  \item \textbf{Codice}: il pattern a memoria, poi laboratorio ed esercizi. Questa è la
        parte più lunga e più importante: il laboratorio è la prova più selettiva.
  \item \textbf{Orale}: 5--10 domande ad alta voce, anche dei moduli passati.
  \item Nei tempi morti (autobus, pause) usa il \textbf{telefono}: quiz, orale e
        \q{Prevedi l'output}. A fine giornata sincronizza i progressi con il PC.
  \item Una volta a settimana, e ogni giorno nell'ultima settimana, fai la
        \textbf{simulazione del quiz} (30 domande, 60 minuti) e il laboratorio della
        simulazione d'esame.
\end{enumerate}
\end{passopasso}

\begin{trappola}[: i punti non sono il voto]
XP, livelli e traguardi servono a darti costanza, non a misurare la preparazione. La
misura vera è la simulazione: punta a superare stabilmente le 22 risposte giuste su 30 e
a scrivere il laboratorio della simulazione in due ore, senza guardare la soluzione.
\end{trappola}

% ---------------------------------------------------------------------
\section*{XP, livelli, missioni e il compagno Bit}
\addcontentsline{toc}{section}{XP, livelli, missioni e il compagno Bit}
% ---------------------------------------------------------------------

Ogni cosa che fai nell'app vale \textbf{punti esperienza} (XP). I punti non si comprano e
non si perdono: sono calcolati dai tuoi progressi, e infatti dopo la sincronizzazione PC e
telefono mostrano sempre lo stesso numero.
\begin{center}\small
\begin{tabular}{lr@{\qquad}lr}
\toprule
\textbf{Attività} & \textbf{XP} & \textbf{Attività} & \textbf{XP} \\
\midrule
risposta giusta al quiz & 10 & esercizio al PC completato & 60 \\
risposta sbagliata (ci hai provato) & 2 & pattern riscritto & 50 \\
domanda d'orale ripassata & 3 & laboratorio a step & 200 \\
ogni livello di una domanda d'orale & 15 & laboratorio della simulazione & 400 \\
output indovinato (telefono) & 25 & simulazione del quiz & 5 a risposta \\
missione del giorno & 40 & simulazione superata & +150 \\
\bottomrule
\end{tabular}
\end{center}
Ogni domanda dà punti solo per i primi cinque tentativi: ripetere all'infinito la stessa
domanda non paga, affrontarne di nuove sì.

\begin{description}
  \item[Livelli.] Il livello $L$ richiede $50 \cdot L \cdot (L-1)$ XP (100, 300, 600,
        1000, \ldots): all'inizio si sale in fretta, poi serve costanza. Ogni fascia ha un
        titolo, da \emph{Hello World} a \emph{Leggenda di Prog2}.
  \item[Missioni.] Ogni giorno tre missioni (una sul quiz, una sull'orale, una sul
        codice), uguali su PC e telefono: 40 XP l'una, 60 in più se le completi tutte.
  \item[Serie.] I giorni consecutivi in cui hai studiato. Non si azzera finché non salti
        un giorno intero.
  \item[Traguardi.] 25 obiettivi da sbloccare: la prima simulazione superata, tutti i
        laboratori, 7 giorni di fila, 30/30 alla simulazione\ldots
  \item[Bit.] Il compagno in pixel art: nasce da un uovo e si evolve con i livelli (Bit,
        Byte, Kilobyte, Megabyte, Gigabyte), ogni volta con forma e accessori nuovi e un
        lampo di evoluzione. Il corpo prende il colore che hai scelto. Respira, sbatte le
        palpebre e salta quando guadagni XP. Si nutre degli XP di oggi rispetto al tuo
        \emph{obiettivo giornaliero} (Impostazioni): se non studi per due giorni ha fame
        e lo stomaco gli brontola, dopo quattro si addormenta (con le Z che fluttuano).
        Ogni tanto ti dà un consiglio sull'esame. Puoi rinominarlo.
\end{description}

% ---------------------------------------------------------------------
\section*{Musica e colori}
\addcontentsline{toc}{section}{Musica e colori}
% ---------------------------------------------------------------------

Il lettore musicale si apre dalla barra laterale (sotto il livello) o da Impostazioni.
Tre sorgenti:
\begin{itemize}
  \item \textbf{Colonna sonora originale}: tre brani composti per l'app e generati dal
        vivo dal computer (\emph{Concentrazione}, lo-fi; \emph{Codice notturno},
        synthwave; \emph{Deep focus}, ambient). Non servono Internet né account.
  \item \textbf{Spotify}: il lettore ufficiale incorporato. Parte con la playlist del
        progetto; incolla il link di una tua playlist (in Spotify: Condividi $\to$ Copia
        link) e premi \emph{Usa}. Per ascoltare i brani interi premi \textbf{Accedi a
        Spotify} e fai l'accesso con email e password (senza accesso Spotify fa sentire
        30 secondi per brano). L'accesso resta valido ai riavvii.
  \item \textbf{YouTube}: per chi non ha Spotify. La playlist suona in sottofondo da una
        piccola miniatura, anche con il lettore chiuso.
\end{itemize}
\texttt{Ctrl+M} toglie e rimette l'audio ovunque ti trovi. In Impostazioni scegli se far
partire la musica all'avvio e se sentire gli effetti sonori di XP e livelli.

Il \textbf{colore} del terminale (verde fosforo all'inizio) si cambia in Impostazioni
$\to$ Aspetto: otto colori pronti o uno a scelta dalla tavolozza. Lo stesso colore vale
sul telefono, se sincronizzi.

% ---------------------------------------------------------------------
\section*{L'Allenatore sul telefono (Android e iPhone)}
\addcontentsline{toc}{section}{L'Allenatore sul telefono (Android e iPhone)}
% ---------------------------------------------------------------------

La versione per telefono è una pagina web che si installa come un'app: niente store,
niente account. Contiene quiz, orale, \textbf{Prevedi l'output} (leggi un programma,
scrivi su carta cosa stampa, poi controlla: è il tipo di domanda \q{qual è l'output?}
del quiz d'esame), gli esercizi in lettura con la soluzione, Bit e la musica. Non
compila: il codice si scrive al PC.

\begin{passopasso}[: installarla]
\begin{enumerate}
  \item Apri dal telefono \url{https://kurtsixxxx.github.io/MyUniversity/} (oppure
        inquadra il codice QR in Impostazioni $\to$ Telefono dell'app per PC).
  \item \textbf{iPhone / iPad}: in \textbf{Safari} tocca \emph{Condividi} (il quadrato
        con la freccia) $\to$ \textbf{Aggiungi alla schermata Home} $\to$ \emph{Aggiungi}.
  \item \textbf{Android}: in \textbf{Chrome} tocca il menu \emph{$\vdots$} $\to$
        \textbf{Installa app} (su alcuni telefoni: \emph{Aggiungi a schermata Home}).
        Se compare il pulsante \emph{Installa l'app} nelle impostazioni dell'app, basta
        quello.
  \item Da quel momento si apre dall'icona, a schermo intero, e funziona anche
        \textbf{senza connessione} (tranne Spotify e YouTube).
\end{enumerate}
\end{passopasso}

\subsection*{Sincronizzare PC e telefono}

I progressi restano sul dispositivo in cui li fai: nessun server, nessun account. Per
unirli si usa un \textbf{codice} di testo che inizia con \texttt{P2:}.
\begin{itemize}
  \item \textbf{Dal PC al telefono}: nell'app per PC, Impostazioni $\to$ Telefono $\to$
        \emph{Copia il codice per il telefono}; mandatelo (per esempio con un messaggio a
        te stesso), poi sul telefono Impostazioni $\to$ \emph{Dal PC al telefono} $\to$
        incolla $\to$ \emph{Unisci}.
  \item \textbf{Dal telefono al PC}: il percorso inverso (\emph{Copia il codice} sul
        telefono, \emph{Unisci i progressi del telefono} sul PC).
\end{itemize}
Le due cronologie vengono \textbf{unite}, non sovrascritte: per ogni domanda vale il
tentativo più recente. Il codice porta con sé anche nome, appello, docente e colore, che
riempiono i campi ancora vuoti dall'altra parte.

% ---------------------------------------------------------------------
\section*{Scorciatoie da tastiera (PC)}
\addcontentsline{toc}{section}{Scorciatoie da tastiera (PC)}
% ---------------------------------------------------------------------

\begin{center}\small
\begin{tabular}{ll@{\qquad}ll}
\toprule
\texttt{Ctrl+1}\ldots\texttt{5} & cambia schermata & \texttt{Ctrl+Invio} / \texttt{F5} & compila ed esegue \\
\texttt{Ctrl+,} & impostazioni & \texttt{Ctrl+S} & salva il codice \\
\texttt{Ctrl+M} & audio sì/no & \texttt{Ctrl+K} & cerca un esercizio \\
\texttt{Ctrl+B} & riduce la barra laterale & \texttt{Ctrl+E} & nasconde l'elenco esercizi \\
\texttt{A}--\texttt{D}, Invio & rispondi e vai avanti (quiz) & \texttt{Ctrl+I} & importa da VS Code \\
Spazio, \texttt{1}--\texttt{3} & mostra e valuta (orale) & \texttt{Ctrl+/} & commenta le righe \\
\bottomrule
\end{tabular}
\end{center}

% ---------------------------------------------------------------------
\section*{Problemi frequenti}
\addcontentsline{toc}{section}{Problemi frequenti}
% ---------------------------------------------------------------------

\begin{description}
  \item[\q{Non trovo il compilatore g++}.] Rifai il passo del \texttt{Path} e
        \textbf{riavvia} l'app (e se non basta il PC): il \texttt{Path} viene letto solo
        all'avvio.
  \item[L'esecuzione dice \q{Permission denied} o non parte.] Spesso è l'antivirus che
        blocca il programma appena compilato scambiandolo per sospetto. Riprova; se
        succede sempre con lo stesso esercizio, cambia un po' il codice (anche solo il
        nome di una variabile) e ricompila.
  \item[\q{Non trovo i contenuti}.] L'app non trova la cartella \texttt{Dispensa}: in
        Impostazioni $\to$ Sistema premi \emph{Scegli la cartella della dispensa}.
  \item[Spotify suona solo 30 secondi.] Manca l'accesso: \emph{Accedi a Spotify} nel
        lettore. Usa email e password: i pulsanti \q{Continua con Google/Apple} aprono il
        browser e l'accesso non arriva all'app.
  \item[YouTube salta un video.] Alcuni autori vietano la riproduzione fuori da YouTube:
        il lettore passa al successivo.
  \item[Dove sono i miei progressi?] Sul PC in \cc|%APPDATA%\Allenatore Prog2| (Impostazioni
        $\to$ I tuoi dati $\to$ Apri la cartella): copiala per fare un backup. Sul telefono
        nella memoria del browser: non cancellare i dati del sito, e ogni tanto copia il
        codice \texttt{P2:} in un posto sicuro.
  \item[Uso Mac o Linux.] L'eseguibile è solo per Windows, ma la versione per telefono
        funziona in qualunque browser, anche sul computer. Chi ha Python può avviare l'app
        dai sorgenti (\cc|python app.py| nella cartella \texttt{Allenatore}, dopo
        \cc|pip install pywebview|).
\end{description}

\begin{consiglio}[: migliorare il kit]
Hai trovato un errore nella dispensa, una domanda ambigua o un esercizio che non compila?
Segnalalo nella pagina del progetto: i sorgenti \LaTeX\ della dispensa e il codice
dell'app sono pubblici, e ogni correzione aiuta chi studierà dopo di te. Per ricompilare
i PDF serve una distribuzione \LaTeX\ (per esempio MiKTeX): doppio clic su
\texttt{compila.bat} nella cartella \texttt{Dispensa}.
\end{consiglio}

\subsection*{Crediti}
Dispensa e Allenatore Prog2 sono un progetto di \textbf{Giulio Casella}, realizzato con
l'aiuto di Claude (Anthropic) e messo a disposizione di tutti gli studenti del corso.
